\section{Discussion}

This pilot extends the 2024 GPT-4 abstract in a narrower but more
auditable direction than originally planned. The main empirical change
is corpus retrieval: the mentor-provided Pushshift-format r/Endo archive
contains far more analyzable material than the 2024 abstract's reported
1{,}986 posts, even after restriction to the overlapping 2018-02-07 to
2022-12-31 window. The model comparison we initially intended could not
be completed under free-tier Gemini limits, so the present claim is not
``two models independently recovered the same themes.'' The defensible
claim is instead that a single contemporary extractor, run over 74
large corpus chunks and checked by two independent judge models, recovers
a stable set of treatment- and care-related themes that can be audited
against source quotes and prior qualitative literature.

The strongest result is within-run recurrence. All 20 final themes were
supported by at least 40 of 74 disjoint chunks, and the highest-support
themes appeared in nearly every slice of the sample. This matters
because a theme that reappears across many independently processed
chunks is less likely to be a one-off artifact of a single prompt
completion. Chunk-support is not a substitute for cross-model
agreement, but it is a useful single-extractor robustness signal and is
more transparent than treating the final reduced theme list as a flat
set of equally supported findings.

The validation stack also clarifies where the current run is weak.
Rotating LLM-as-judge scores were high, especially for faithfulness and
clinical relevance, but quote grounding was only 35\% by exact-match
and BM25 retrieval. Some of that gap is expected because embedding-based
semantic grounding was disabled, and because the reduce step may carry
lightly edited quotes forward. Still, the low lexical grounding rate is
important: the judge scores should not be read as proof that every
supporting quote is audit-ready. In a revision intended for submission,
the quote fields should either be re-grounded with semantic retrieval or
manually replaced with verified paraphrases before publication.

The Young~2015 anchor results separate broad convergence from
fine-grained online-community topics. Nine of 20 final themes aligned
lexically with the prior qualitative synthesis, including diagnostic
delay, medical dismissal, chronic pain, sexual dysfunction, fertility
concerns, psychological impact, and peer support. Themes that did not
align, such as laparoscopic excision, GnRH-based management,
post-operative recovery, recurrence after surgery, pelvic floor
therapy, dietary changes, bowel involvement, hysterectomy, specialist
access, and adenomyosis comorbidity, are not automatically novel
findings. They are more precise treatment-management topics than the
Young~2015 theme set was designed to encode. Their value is practical:
they suggest concrete issues clinicians and researchers can ask about
when discussing endometriosis care with patients.

The most clinically useful themes are therefore not necessarily the
highest composite-reliability themes. Diagnostic delay and medical
dismissal are both highly reliable and already well represented in the
literature. In contrast, specialist access, pelvic floor dysfunction,
post-operative recovery, and adenomyosis comorbidity score lower
because they do not match the Young anchor or have weaker lexical quote
grounding, yet they still appear across many chunks and may be
important prompts for future qualitative review. The composite score
should be interpreted as a triage aid for human review, not a final
truth label.

This study is best understood as a reproducibility and methods pilot.
It shows that a low-budget, single-extractor workflow can produce a
traceable theme table, preserve intermediate chunk-level outputs, and
attach multiple reliability signals to every final theme. It also shows
where such a workflow falls short: cross-model agreement, prompt
paraphrase stability, high-temperature stability, and human
patient/clinician review are still absent. A funded follow-up should
run a second extractor on the same frozen sample, repeat the prompt-B
and temperature-stability sweeps, and calibrate Sonnet/Opus judge
scores against human ratings.

Several limitations bound the strength of the findings. The analysis is
restricted to r/Endo because r/endometriosis was not present in the
mentor-provided archive, and the archive ends on 2022-12-31, leaving the
2023-01 to 2024-10-31 portion of the published window unanalyzed. Reddit
users are self-selected, English-speaking, and more digitally engaged
than the broader endometriosis population. The original 2024 GPT-4 run
could not be reproduced without an OpenAI credential. Finally, although
we now preserve the subagent protocol for reruns, the exact interactive
subagent messages from the completed run were not captured as a
first-class artifact. Future runs should log those prompts before any
subagent work begins.
